\section{Разработка системы распознавания языка}
\label{sec:development}

Система разворачивается в локальной вычислительной сети средствами Docker Compose и состоит из трёх сервисов: PostgreSQL~16, сервер FastAPI на Python~3.11 и одностраничный интерфейс на React за nginx.

Серверная часть использует FastAPI и Uvicorn, доступ к базе \mdash \texttt{psycopg2}. Из PDF текст извлекается библиотекой \texttt{pypdf}. Русская лемматизация выполняется \texttt{pymorphy3}. Нейросеть \mdash \texttt{MLPClassifier} из \texttt{scikit-learn}. Графики метрик строит \texttt{matplotlib} и сохраняет в каталог отчёта \texttt{png/}. Клиент \mdash React~18, Vite и \texttt{react-router-dom}. Требование варианта о формате документов закрыто ограничением загрузки: на вход принимается только PDF.

Структура системы показана на рисунке~\ref{fig:architecture}. Модуль \texttt{main.py} объявляет REST-эндпоинты. \texttt{text\_processing.py} выполняет токенизацию и лемматизацию. \texttt{language\_model.py} реализует три метода и расстояние ОоП. \texttt{document\_loader.py} читает PDF. \texttt{indexer.py} сохраняет документы и помечает источник: эталон из \texttt{qrels.json} или загруженный файл. \texttt{evaluator.py} считает Accuracy, точность, полноту, F-меру и время. \texttt{bootstrap.py} при старте загружает обучающий корпус, обучает модели и индексирует эталонные фразы. \texttt{migrate.py} применяет SQL-миграции.

\begin{figure}[H]
	\centering
	\includegraphics[width=1.0\textwidth]{architecture.png}
	\caption{Структура системы распознавания языка.}
	\label{fig:architecture}
\end{figure}

Последовательность инициализации при старте сервера показана на рисунке~\ref{fig:seq-train}: корпус читается в \texttt{training\_texts}, строятся ПОЯ и MLP, эталонные фразы из \texttt{qrels.json} попадают в \texttt{eval\_queries} и в список документов как \texttt{eval\_*.txt}.

\begin{figure}[H]
	\centering
	\includegraphics[width=1.0\textwidth]{sequence_train.png}
	\caption{Диаграмма последовательности обучения моделей и индексации эталонов.}
	\label{fig:seq-train}
\end{figure}

Загрузка PDF и распознавание показаны на рисунке~\ref{fig:seq-search}. Пользователь выбирает файл на странице распознавания; nginx передаёт запрос на \texttt{POST /api/recognize}; сервер извлекает текст, прогоняет три метода, записывает документ и возвращает язык, уверенность, ОоП и время.

\begin{figure}[H]
	\centering
	\includegraphics[width=1.0\textwidth]{sequence_search.png}
	\caption{Диаграмма последовательности распознавания PDF-документа.}
	\label{fig:seq-search}
\end{figure}

Таблицы базы перечислены в таблице~\ref{tab:db-tables}, схема \mdash на рисунке~\ref{fig:db}. Обучающие тексты хранятся в \texttt{training\_texts} и не попадают в пользовательский список. В \texttt{documents} лежат эталонные фразы (файлы \texttt{eval\_*.txt}) и загруженные PDF.

\begin{table}[H]
	\centering
	\begin{tabular}{|p{4.2cm}|p{8.8cm}|}
		\hline
		\textbf{Таблица} & \textbf{Назначение} \\
		\hline
		\texttt{documents} & Заголовок, текст, язык, результат, источник \\
		\hline
		\texttt{training\_texts} & Обучающие тексты по языкам \\
		\hline
		\texttt{language\_profiles} & Лексемы и частоты ПОЯ по методам \\
		\hline
		\texttt{recognition\_results} & Результат каждого метода по документу \\
		\hline
		\texttt{eval\_queries} & Эталонные фразы и метка языка \\
		\hline
		\texttt{eval\_runs} & Сводка прогона оценки \\
		\hline
	\end{tabular}
	\caption{Таблицы базы данных системы распознавания языка.}
	\label{tab:db-tables}
\end{table}

\begin{figure}[H]
	\centering
	\includegraphics[width=1.0\textwidth]{db_schema.png}
	\caption{Схема базы данных.}
	\label{fig:db}
\end{figure}

Интерфейс содержит страницы распознавания, списка документов с ссылкой на карточку, оценки качества и справки. Список документов можно сохранить в JSON и напечатать из браузера. Развёртывание: \texttt{docker compose up -{}-build} из каталога программной части. Интерфейс доступен по адресу \texttt{http://localhost:8080}.

Проверка \texttt{scenarios.py} покрывает лемматизацию русского текста, сохранение немецких умлаутов, оба словесных метода, ОоП по рангам и обучение нейросети на короткой выборке. Все десять проверок прошли.

Оценка качества выполняется на десяти эталонных фразах (пять русских и пять немецких). Каждый метод классифицирует фразу, после чего считаются Accuracy, макросредние точность, полнота и F-мера, а также среднее время. Результаты \mdash в таблице~\ref{tab:metrics}. На этой выборке все три метода дали полную точность: русский и немецкий различаются алфавитом, и профили почти не пересекаются.

\begin{table}[H]
	\centering
	\begin{tabular}{|l|c|c|c|c|c|}
		\hline
		\textbf{Метод} & \textbf{Acc} & \textbf{P} & \textbf{R} & \textbf{F1} & \textbf{мс} \\
		\hline
		Коротких слов & 1,000 & 1,000 & 1,000 & 1,000 & 0,26 \\
		\hline
		Частотных слов & 1,000 & 1,000 & 1,000 & 1,000 & 0,30 \\
		\hline
		Нейросетевой & 1,000 & 1,000 & 1,000 & 1,000 & 30,17 \\
		\hline
	\end{tabular}
	\caption{Метрики на эталонной выборке из десяти фраз.}
	\label{tab:metrics}
\end{table}

На длинном PDF картина по времени другая: метод коротких слов обходит весь поток коротких лексем и широкий ПОЯ, частотный ограничен топ-50, нейросеть строит N-граммы по всему тексту. На демонстрационном документе получены величины порядка 355~мс, 58~мс и 403~мс соответственно. Сравнение методов по качеству и по времени показано на рисунках~\ref{fig:method-cmp} и~\ref{fig:time-cmp}.

\begin{figure}[H]
	\centering
	\includegraphics[width=1.0\textwidth]{method_comparison.png}
	\caption{Сравнение методов по Accuracy, точности, полноте и F-мере.}
	\label{fig:method-cmp}
\end{figure}

\begin{figure}[H]
	\centering
	\includegraphics[width=0.85\textwidth]{time_comparison.png}
	\caption{Среднее время методов на эталонных фразах.}
	\label{fig:time-cmp}
\end{figure}

Для ручной проверки подготовлены четыре одностраничных PDF в каталоге \texttt{demo-pdf}: два русских (\texttt{ru\_priem.pdf}, \texttt{ru\_operaciya.pdf}) и два немецких (\texttt{de\_aufnahme.pdf}, \texttt{de\_operation.pdf}). Их можно загрузить на странице распознавания.
